\begin{proof}
\begin{enumerate}[label=\arabic*.]
\item Write the finite generators of $W$ as $S_1,S_2$ and the
extension generators as
\[
 A_i=\begin{pmatrix}1&a_i\\0&S_i\end{pmatrix}.
\]
If $S_i$ has order $m_i$, retaining that finite order is equivalent
to $a_i(I+S_i+\cdots+S_i^{m_i-1})=0$.  Averaging identifies this
condition with $a_i\in\operatorname{row}(S_i-I)$, a space of
dimension four.  A change of splitting by $b\in W^*$ sends the
two rows to $a_i+b(S_i-I)$.  This coboundary map is injective
because $W$ has no invariant covectors.  Thus the extension
space subject to the two finite conditions has dimension $8-5=3$.

\item Put $U=S_1S_2=I+N$, so $N^2=0$ and $\operatorname{rank}N=2$.
The product $A_1A_2$ has upper row $c=a_1S_2+a_2$.
Changing the splitting sends $c$ to $c+bN$.  Hence cusp restriction
is the explicit map
\[
 [(a_1,a_2)]\longmapsto[c]\in W^*/\operatorname{row}N.
\]
It is an isomorphism.  Indeed, a class in its kernel can be
represented with $c=0$, so $a_2=-a_1S_2$.  The finite conditions
then say
\[
 a_1\in\operatorname{row}(S_1-I)
              \cap\operatorname{row}(S_2-I).
\]
This intersection has dimension three: the two four-dimensional
row spaces span $W^*$ because $W$ has no common fixed vector.
The rows $b$ satisfying $bN=0$ also form a three-dimensional
space.  Their map $b\mapsto b(S_1-I)$ lands in the displayed
intersection, since $bU=b$ gives
$b(S_1-I)S_2=b-bS_2$.  It is injective, because
$bS_1=b$ and $bU=b$ together give $bS_2=b$, forcing $b=0$.
It is therefore onto.  Every pair with $c=0$ is a coboundary,
proving injectivity of cusp restriction.  Both its source and
target have dimension three, so it is an isomorphism.

\item The square of the lifted product minus identity vanishes
exactly when $cN=0$.  Under this restriction isomorphism the
square-zero extensions form
\[
 \frac{\operatorname{ann}(\operatorname{im}N)}
      {\operatorname{row}N},
 \tag{F.32}\label{maps:eq:F.32}
\]
of dimension $3-2=1$.  The lifted nilpotent has the rows of $N$
together with the row $c$, so its rank is three exactly when
$c\notin\operatorname{row}N$.  These are precisely the nonzero
classes in \eqref{maps:eq:F.32}.  Rescaling the invariant line
rescales this one-dimensional rational space and is transitive
on its nonzero classes.  This proves the asserted existence and
uniqueness with the stated finite orders and maximal cusp.
\end{enumerate}
\end{proof}
